\documentclass[aps,pra,twocolumn,superscriptaddress,longbibliography]{revtex4-2}
\usepackage{amsmath,amssymb,graphicx,xcolor,bm}

\begin{document}

\title{The Dimensionless Vacuum: Recovering the Schr\"odinger Equation\\from Scale-Invariant Graph Topology}
\thanks{Paper 7 in the Geometric Vacuum series}

\author{J.~Loutey}
\affiliation{Independent Researcher, Kent, Washington}

\date{March 10, 2026}

\begin{abstract}
Fock showed in 1935 that the hydrogen Schr\"odinger equation is mathematically equivalent to the Laplace-Beltrami operator on the unit three-sphere $S^3$ via stereographic projection. \textbf{[MEASURED]} We report what the GeoVac discrete graph Laplacian is measured to do through $n_{\max} = 30$---its top eigenvalue saturates the bipartite bound $\lambda_{\max} \to 2 d_{\max} = 8$, ---and \textbf{[OBSERVATION]} we read this as a numerical indication that the lattice converges to this same operator in the continuum limit; the operator convergence itself is not established here (2026-09-03 correction: an earlier version stated it as demonstrated). We verify the conformal chain from the continuum operator onward by symbolic computation. The discrete paraboloid lattice---whose nodes encode quantum numbers $(n, l, m)$ and whose edges encode transition amplitudes---is read as possessing a continuum limit conformally equivalent to the unit $S^3$ (a reading with numerical indications, not an established convergence; item 1 below). \textbf{[SYMBOLIC PROOF]} The continuum $S^3$ topology (the limit of the positive-semidefinite graph Laplacian) has pure integer eigenvalues $\lambda_n = -(n^2 - 1)$ with no intrinsic physical scale. \textbf{[SYMBOLIC PROOF]} The familiar Rydberg energy levels $E_n = -1/(2n^2)$ emerge only upon imposing an energy-shell constraint $p_0^2 = -2E$ that acts as a ``stereographic focal length,'' projecting the curved topological information into flat $\mathbb{R}^3$ observational coordinates. \textbf{[SYMBOLIC PROOF]} We verify these claims through a suite of 18 symbolic checks (an algebraic audit of the transformation chain, several of them limiting-case or definitional) confirming the algebraic geometry of the projection, the operator equivalence under conformal rescaling, and the eigenvalue-to-energy mapping.
\end{abstract}

\maketitle


%% ========================================================================
\section{Introduction}
\label{sec:intro}
%% ========================================================================

The Schr\"odinger equation is the foundation of non-relativistic quantum mechanics. In atomic units ($\hbar = m_e = e = 1$), the time-independent equation for hydrogen reads
\begin{equation}
\left( -\frac{1}{2}\nabla^2 - \frac{1}{r} \right) \psi(\mathbf{r}) = E\,\psi(\mathbf{r}).
\label{eq:schrodinger}
\end{equation}
This equation presupposes a continuous, flat Euclidean space $\mathbb{R}^3$ and a radial Coulomb potential $V(r) = -1/r$ with explicit physical dimensions. These assumptions are so deeply embedded in the formalism that they are rarely questioned: continuous space provides the arena, the Coulomb potential provides the dynamics, and the quantized energy levels $E_n = -1/(2n^2)$ emerge as boundary-condition eigenvalues.

Yet a remarkable observation, first made by Fock in 1935 \cite{fock1935}, suggests a fundamentally different perspective. Fock showed that when hydrogen is analyzed in momentum space, the Schr\"odinger equation acquires an unexpected geometric structure. The momentum-space integral equation, with its $1/|\mathbf{p} - \mathbf{p}'|^2$ kernel (the Fourier transform of the Coulomb potential), can be mapped by stereographic projection onto the three-sphere $S^3$. On $S^3$, the Coulomb problem reduces to a \textit{free particle}: the eigenfunctions become hyperspherical harmonics, and the spectrum is determined by the topology of the sphere alone. The $1/r$ potential is not fundamental---it is a projection artifact of curved geometry viewed in flat coordinates.

This observation was developed further by Bargmann \cite{bargmann1936}, who connected the hidden $SO(4)$ symmetry of hydrogen to the isometry group of $S^3$, and by Barut and Kleinert \cite{barut1967}, who extended the symmetry to the full conformal group $SO(4,2)$. The mathematical structure has been thoroughly studied in the context of dynamical symmetry groups \cite{bander1966,biedenharn1981}. However, these results are typically presented as mathematical curiosities---elegant reformulations of a solved problem---rather than as evidence for a deeper physical ontology.

In this paper, we take the opposite stance. We argue that the $S^3$ representation is not merely a convenient mathematical trick, but reveals deep structural content that can be exploited computationally. Our starting point is the GeoVac framework \cite{loutey_paper0,loutey_paper1}, which models the hydrogen atom as a discrete graph: nodes represent quantum states $|n, l, m\rangle$, edges represent physical transitions, and the Hamiltonian is a sparse matrix encoding connectivity. This discrete lattice is \textit{intrinsically dimensionless}---it is a pure combinatorial topology with no inherent length, time, or energy scales.

The central results of this paper are one empirical convergence statement and two symbolic ones:
\begin{enumerate}
\item \textbf{[MEASURED]} The top Laplacian eigenvalue of the discrete graph saturates its bipartite bound, $\lambda_{\max} \to 2 d_{\max} = 8$, through $n_{\max} = 30$.  \textbf{[OBSERVATION]} We read this as a numerical indication of a continuum limit conformally equivalent to the \textit{unit} $S^3$;\ the operator convergence and the identification of the limit manifold are \emph{not} established here (see New Contributions item~1).  \emph{Corrected 2026-09-03/04:} this item previously asserted both as \textbf{[MEASURED]}, contradicting four other loci in this paper.
\item \textbf{[SYMBOLIC PROOF]} The Laplace-Beltrami operator on this unit sphere produces pure integer eigenvalues $\lambda_n = -(n^2 - 1)$, carrying no physical dimensions.
\item \textbf{[SYMBOLIC PROOF]} The dimensionful Rydberg spectrum $E_n = -1/(2n^2)$ is introduced exclusively through the energy-shell constraint $p_0^2 = -2E$, which serves as a projection parameter mapping $S^3$ coordinates to flat $\mathbb{R}^3$ coordinates.
\end{enumerate}

These claims are supported by a suite of 18 symbolic checks, implemented in \texttt{sympy} and verified computationally, that confirm every algebraic step of the conformal transformation chain (several are limiting-case or definitional substitutions rather than logically independent theorems).

The structure of this paper is as follows. Section~\ref{sec:graph} describes the dimensionless graph topology and its relationship to $S^3$. Section~\ref{sec:conformal} details the stereographic projection and the role of the conformal factor. Section~\ref{sec:operator} demonstrates the operator equivalence between the curved Laplacian and the flat Schr\"odinger equation. Section~\ref{sec:vee_s3} derives the two-electron repulsion integral as a density-overlap integral on $S^3$. Section~\ref{sec:discussion} discusses the physical and philosophical implications.


%% ========================================================================
\section{What is New vs.\ What is Known}
\label{sec:attribution}
%% ========================================================================

The $S^3$ structure of the hydrogen atom has a distinguished history, and we emphasize what is established mathematics versus what is new in this work.

\subsection{Known Results (Prior Work)}

The following results are well-established in the mathematical physics literature:

\begin{itemize}
\item \textit{Fock (1935)} \cite{fock1935}: The hydrogen atom in momentum space maps to $S^3$ via stereographic projection. The Coulomb problem becomes a free particle on the sphere, with the $1/r$ potential emerging from the chordal distance identity.

\item \textit{Bargmann (1936)} \cite{bargmann1936}: The hidden $SO(4)$ symmetry of hydrogen corresponds to the isometry group of $S^3$. The integer spectrum $\lambda_n = -(n^2 - 1)$ follows from the representation theory of $SO(4)$.

\item \textit{Barut and Kleinert (1967)} \cite{barut1967,barut1967b}: The full conformal group $SO(4,2)$ acts as the dynamical symmetry group of hydrogen, connecting $S^3$ to the non-compact spectrum-generating algebra.
\end{itemize}

The stereographic projection formulas (Eqs.~\ref{eq:stereo_spatial}--\ref{eq:stereo_fourth}), the unit sphere constraint (Eq.~\ref{eq:unit_sphere}), the chordal distance identity (Eq.~\ref{eq:chordal}), the conformal Laplacian relation (Eq.~\ref{eq:conformal_laplacian}), and the integer eigenvalue spectrum (Eq.~\ref{eq:s3_spectrum}) are all known results. Our symbolic proofs in the Appendix verify these identities; they do not claim to derive them for the first time.

\subsection{New Contributions of This Work}

The following are the specific new contributions:

\begin{enumerate}
\item \textbf{Discrete graph convergence.} \textbf{[OBSERVATION]} We read the GeoVac discrete paraboloid lattice---with nodes $(n, l, m)$ and edges generated by $L_\pm, T_\pm$ operators---as a discretization whose continuum limit is the Laplace-Beltrami operator on unit $S^3$, connecting a concrete computational framework (sparse graph eigenvalue problems) to the known continuous structure. What is \textbf{[MEASURED]} through $n_{\max} = 30$ (\texttt{tests/test\_paper7\_graph\_convergence.py}; \texttt{tests/test\_trunk\_qa\_splitting.py}, Paper~1) is narrower than that reading: (i) \textbf{[INTERNAL THEOREM]} the top Laplacian eigenvalue of the production lattice saturates the bipartite bound, $\lambda_{\max} \to 2 d_{\max} = 8$ --- the $\ell$-block grid structure makes the spectrum closed form, so this is proved rather than sampled --- (6.618 at $n_{\max} = 5$, 7.954 at 30), at the closed-form rate $(C + o(1))/n_{\max}^2$ with $C = \tfrac{\pi^2}{4}(2+2^{1/3})^2(1+2^{-2/3}) = 42.7397\ldots$ that the $\ell$-block grid structure supplies (the approach is monotone from below, so the $n_{\max} = 320$ sample $42.6$ printed here before 2026-09-03 understated it) (Paper~0~\S VI:\ each $\ell$-block is the grid graph $P_{n_{\max}-\ell} \times P_{2\ell+1}$, so the spectrum is closed form), so that with the matched scale $\kappa = -1/16$ the extremal eigenvalue $\kappa\lambda_{\max}$ approaches $-1/2$ \emph{by construction}---the content is the saturation, not the target; (ii) the $2s/2p$ node-amplitude proxy on the binary lattice decays, oscillating, to $0.39\%$ at $n_{\max} = 30$ (a proxy, not a spectral gap:\ the $\ell$-blocks are disconnected and the mode selection is a near-tie --- Paper~1 \S III conditioning caveat). Two measured facts bound the reading: the lattice's edges never change $l$, so $L$ splits into $n_{\max}$ blocks (one per $l$) with an $n_{\max}$-dimensional kernel of block constants, and the $\lambda_{\max}$ mode lives in a mid-$l$ block with negligible weight on the $1s$ node (the $l = 0$ block alone saturates only $\lambda_{\max} \to 4$); and the spectrum of $\kappa L$ is confined to $[-1/2, 0]$, so no level-by-level Rydberg ladder is reproduced. The operator convergence $L \to \Delta_{S^3}$ and the identification of the limit manifold are not established here (coverage gap; 2026-09-03 correction of a 2026-09-02 backing note that presented the matched target $-1/2$ as the measured quantity, reversing Paper~0's own ``matched, not derived'' scope).

\item \textbf{Explicit symbolic verification.} \textbf{[SYMBOLIC PROOF]} The 18 symbolic proofs (Appendix~\ref{app:tests}) provide a machine-verified algebraic audit of the conformal transformation chain (ten of the eighteen checks are definitional, limiting-case or non-discriminating rather than logically independent theorems). While the underlying identities are known, this systematic verification in \texttt{sympy} is new and provides a reproducible computational check.

\item \textbf{Scale-invariance framing.} We articulate the ``dimensionless vacuum'' perspective: the observation that the conformal factor $\Omega$ always produces the \textit{unit} $S^3$ regardless of $p_0$, and that all physical dimensions enter exclusively through the energy-shell constraint. \textbf{[OBSERVATION]} While implicit in Fock's original treatment, this framing---and its implications for the GeoVac computational approach---has not, to our knowledge, been stated explicitly in the literature.

\item \textbf{Computational methodology.} \textbf{[MEASURED]} The practical consequence: replacing continuous integration with sparse graph Laplacian eigenvalue problems whose operator size is near-linear in $V$ (adjacency nonzeros scale as $V^{1.07}$; the timing is guarded sub-quadratic, Paper~1), justified by the mathematical equivalence established here.
\end{enumerate}


%% ========================================================================
\section{The Dimensionless Graph}
\label{sec:graph}
%% ========================================================================

\subsection{The Paraboloid Lattice}

The GeoVac framework represents the single-electron hydrogen atom as a finite directed graph $\mathcal{G} = (V, E)$. The vertex set $V$ consists of all quantum states $|n, l, m\rangle$ with $1 \leq n \leq n_{\max}$, $0 \leq l < n$, and $-l \leq m \leq l$, giving $|V| = \sum_{n=1}^{n_{\max}} n^2$ nodes. The edge set $E$ is generated by four transition operators:
\begin{align}
L_\pm |n, l, m\rangle &\propto |n, l, m \pm 1\rangle, \label{eq:angular}\\
T_\pm |n, l, m\rangle &\propto |n \pm 1, l, m\rangle. \label{eq:radial}
\end{align}
These operators form the algebraic skeleton of the $SU(2) \otimes SU(1,1)$ dynamical symmetry group, which is the maximal compact$\,\times\,$non-compact factorization of the hydrogen conformal group $SO(4,2)$ \cite{barut1967}.

The graph Laplacian is defined as
\begin{equation}
L = D - A,
\label{eq:laplacian}
\end{equation}
where $A_{ij}$ is the weighted adjacency matrix constructed from the transition operators and $D$ is the diagonal degree matrix with entries $D_{ii} = \sum_j |A_{ij}|$. The single-particle Hamiltonian is
\begin{equation}
H = \kappa\, L = \kappa(D - A),
\label{eq:hamiltonian}
\end{equation}
where $\kappa = -1/16$ is the kinetic scale factor---the constant that maps the dimensionless graph eigenvalues to the physical hydrogen spectrum (the scale matching the binary graph's $\lambda_{\max} \to 8$ to the ground-state energy $-\tfrac12$).  \textbf{[OBSERVATION]} This matching value \emph{coincides} with a single conformal-geometry value, $1/16$, namely the inverse Fock Jacobian $1/\Omega^4(0)$ with $\Omega(0) = 2$ (Section~\ref{sec:discussion}; equivalently one quarter of the $s$-wave coupling $c^2(n,0) = 1/4$ --- corrected 2026-09-03 from a printed ``$c^2(n,0) = (1/4)^2$'', the Chebyshev amplitude being $1/2$) in the dimensionless $p_0 = 1$ convention used throughout this section (dimensionfully $\Omega(0) = 2/p_0$).  We record this as an \emph{observation}: the matching $\kappa$ and the geometric $1/16$ agree numerically, but are computed from disjoint inputs and no mechanism is known that derives one from the other.

\subsection{Intrinsic Dimensionlessness}

A crucial observation about the graph $\mathcal{G}$ is that it carries \textit{no intrinsic physical scale}. The quantum numbers $(n, l, m)$ are pure integers. The edge weights $w_{ij} \propto 1/(n_i n_j)$ are dimensionless ratios. The graph Laplacian $L$ is a dimensionless matrix. The only dimensionful quantity in Eq.~\eqref{eq:hamiltonian} is the kinetic scale factor $\kappa$, which converts the dimensionless graph eigenvalues into physical energy units. But $\kappa$ is an \textit{external} parameter imposed at the measurement stage---it is not encoded in the graph topology itself.

This dimensionlessness is not accidental. It is a direct consequence of the graph's structure as a discrete analog of the unit three-sphere. To see this, we examine the continuum limit.

\subsection{The Continuum Limit: Unit $S^3$}

The graph $\mathcal{G}$ is a finite lattice truncation of a manifold. \textbf{[OBSERVATION]} As $n_{\max} \to \infty$ we read the discrete graph Laplacian $L$ as converging to the continuous Laplace-Beltrami operator $\Delta_{S^3}$ on a three-dimensional manifold; the reading rests on the one spectral indication that survives scrutiny (the spectral bound that saturates), not on an operator-norm convergence statement, which is not established in this paper, and no longer on the $s/p$ proxy, whose size tracks the residue of $n_{\max}$ modulo $3$ (the QA note in this section;\ withdrawn as evidence 2026-09-03) (New Contributions, item 1).

At finite $n_{\max}$, the graph Laplacian exhibits discretization artifacts: the $2s$/$2p$ node-amplitude proxy reaches $37\%$ at $n_{\max} = 5$ on the binary production lattice (the splitting is strongly non-monotone: $12.7\%$, $0.65\%$, $15.6\%$, $5.1\%$ and $1.7\%$ at $n_{\max} = 6, 7, 8, 9, 10$, then $4.1\%$ at $20$ and $0.39\%$ at $30$; measured 2026-09-03), due to differential connectivity between pole ($l = 0$) and equatorial ($l > 0$) nodes \cite{loutey_paper1}. However, this splitting oscillates and decays systematically:
\begin{equation}
\Delta E_{2s\text{-}2p}(n_{\max}) \to 0 \quad \text{as} \quad n_{\max} \to \infty,
\end{equation}
recovering the exact Coulomb degeneracy.  \textbf{[INTERNAL THEOREM]} The rate is closed form and is governed by the residue of $n_{\max}$ modulo $3$, not by proximity to a continuum limit:\ $\lambda_{2s} = 3$ exactly when $n_{\max} \equiv 0 \ (\mathrm{mod}\ 3)$, where the ratio is $\bigl(2-2\cos\frac{\pi}{n_{\max}-1}\bigr)/3 = O(n_{\max}^{-2})$, while off that class it is $O(n_{\max}^{-1})$ and three to six times larger at neighbouring cutoffs ($1.65\%$ at $28$, $2.58\%$ at $29$).  \emph{The ``spectral aliasing on a compact manifold'' reading printed here before 2026-09-03 is withdrawn:} the $\ell$-blocks are disconnected, so no degeneracy is present to alias (Paper~1 \S III).  \textbf{This leg accordingly carries no evidential weight for the continuum reading}, which now rests on the $\lambda_{\max}$ saturation alone. \textbf{[MEASURED]} At $n_{\max} = 30$, the splitting has decayed to well under 1\%, continuing toward the exact Coulomb degeneracy.

The exact degeneracy \emph{of the continuum operator} is a hallmark of $SO(4)$ symmetry---the isometry group of $S^3$.  \textbf{[MEASURED]} The graph's finite-$n_{\max}$ proxy for it is not:\ its value tracks the residue of $n_{\max}$ modulo $3$ (Paper~1 \S III, 2026-09-03), so its decay is not a recovery of that symmetry. \textbf{[OBSERVATION]} A continuum limit carrying $SO(4)$ isometry is consistent with this, and $S^3$ is its natural candidate (up to discrete quotients); the identification is a reading, not a measurement---no test in the repository establishes the limit manifold, and the production graph is disconnected into $n_{\max}$ components (one per $l$), a structure the reading must eventually account for.

The critical question is: what is the \textit{radius} of this sphere? As we shall demonstrate in Section~\ref{sec:conformal}, the answer is unambiguous: the intrinsic geometry is always that of the \textit{unit} $S^3$, regardless of any physical parameters. This is the mathematical content of scale invariance.


%% ========================================================================
\section{The Conformal Lens: Stereographic Projection}
\label{sec:conformal}
%% ========================================================================

\subsection{Fock's Mapping}

Fock's 1935 insight was to map three-dimensional flat momentum space $\mathbb{R}^3$ onto the unit three-sphere $S^3 \subset \mathbb{R}^4$ via stereographic projection \cite{fock1935}. Given a momentum vector $\mathbf{p} = (p_1, p_2, p_3)$ and an energy-shell parameter $p_0$ (defined below), the projection is:
\begin{align}
n_i &= \frac{2p_0\, p_i}{p^2 + p_0^2}, \quad i = 1, 2, 3, \label{eq:stereo_spatial}\\
n_4 &= \frac{p^2 - p_0^2}{p^2 + p_0^2}, \label{eq:stereo_fourth}
\end{align}
where $p^2 = p_1^2 + p_2^2 + p_3^2$. This map sends the origin $\mathbf{p} = 0$ to the south pole $(0, 0, 0, -1)$ and the point at infinity $|\mathbf{p}| \to \infty$ to the north pole $(0, 0, 0, +1)$, thereby compactifying $\mathbb{R}^3$ into $S^3$.

\subsection{The Unit Sphere Constraint}

A direct algebraic computation confirms that the image lies on the unit sphere:
\begin{equation}
\sum_{i=1}^{4} n_i^2 = \frac{4p_0^2 p^2 + (p^2 - p_0^2)^2}{(p^2 + p_0^2)^2} = 1.
\label{eq:unit_sphere}
\end{equation}
This identity holds for \textit{all} values of $p_0 > 0$. The parameter $p_0$ does not change the target manifold; it only changes which point in $\mathbb{R}^3$ maps to the south pole. \textbf{[SYMBOLIC PROOF]} This is our first key result:

\medskip
\noindent\textit{The stereographic projection always produces the unit $S^3$, regardless of $p_0$. The intrinsic geometry of the target manifold has no physical scale.}
\medskip

To make this explicit, define the conformal factor
\begin{equation}
\Omega(\mathbf{p}) = \frac{2p_0}{p^2 + p_0^2}.
\label{eq:conformal_factor}
\end{equation}
The spatial components of the projection satisfy the identity
\begin{equation}
n_i = \Omega \cdot p_i, \quad i = 1, 2, 3,
\label{eq:ni_omega}
\end{equation}
confirming that $\Omega$ is the local scale factor of the map. Under the substitution $u = p/p_0$ (dimensionless momentum), the conformal factor becomes
\begin{equation}
\Omega = \frac{2}{p_0(1 + u^2)},
\end{equation}
and the induced metric on $S^3$, expressed in flat stereographic coordinates, takes the form
\begin{equation}
\mathrm{d}s^2_{S^3} = \Omega^2\, \mathrm{d}p^2 = \frac{4}{(1 + u^2)^2}\, \mathrm{d}u^2.
\label{eq:unit_metric}
\end{equation}
This is the standard round metric on the unit $S^3$ in stereographic coordinates. \textbf{[SYMBOLIC PROOF]} The factors of $p_0$ cancel identically, confirming that the intrinsic geometry is scale-invariant.

\subsection{The Chordal Distance Identity}

The conformal structure of the projection is captured by a fundamental distance identity. For two momentum vectors $\mathbf{p}$ and $\mathbf{q}$ with images $\mathbf{n}(\mathbf{p})$ and $\mathbf{n}(\mathbf{q})$ on $S^3$:
\begin{equation}
|\mathbf{n}(\mathbf{p}) - \mathbf{n}(\mathbf{q})|^2 = \Omega(\mathbf{p})\;\Omega(\mathbf{q})\; |\mathbf{p} - \mathbf{q}|^2.
\label{eq:chordal}
\end{equation}
The left side is the chordal (Euclidean 4D) distance on $S^3$; the right side is the flat Euclidean distance in $\mathbb{R}^3$, weighted by conformal factors at both endpoints.

This identity is the algebraic engine behind the emergence of the Coulomb potential. The momentum-space Schr\"odinger equation contains the kernel $1/|\mathbf{p} - \mathbf{p}'|^2$, which is the Fourier transform of the $1/r$ Coulomb potential. Under Fock's projection, Eq.~\eqref{eq:chordal} converts this kernel into $1/|\mathbf{n} - \mathbf{n}'|^2$ on $S^3$ (up to conformal weights), transforming the interacting Coulomb problem into a \textit{free-particle} problem on the sphere \cite{fock1935}.

\subsection{Physical Interpretation of $p_0$}

What is $p_0$? In Fock's original treatment, it is defined by the energy-shell constraint:
\begin{equation}
p_0^2 = -2m E,
\label{eq:energy_shell}
\end{equation}
where $E < 0$ for bound states (we set $m = 1$ in atomic units). For the $n$-th energy level of hydrogen, $E_n = -1/(2n^2)$, so $p_0 = 1/n$.

Crucially, $p_0$ does \textit{not} determine the curvature of the target manifold---the unit sphere constraint Eq.~\eqref{eq:unit_sphere} holds for all $p_0$. Instead, $p_0$ acts as a \textbf{stereographic focal length}: it determines how the flat momentum coordinates $\mathbf{p}$ are laid out on the fixed unit sphere. Different values of $p_0$ correspond to different stereographic charts of the \textit{same} underlying manifold, analogous to choosing different focal lengths for a camera lens viewing the same scene.

The conformal factor $\Omega$ encodes this projection:
\begin{itemize}
\item At $\mathbf{p} = 0$: $\Omega = 2/p_0$ (maximum magnification at the south pole).
\item As $|\mathbf{p}| \to \infty$: $\Omega \to 0$ (the north pole is infinitely compressed).
\end{itemize}
The volume element transforms as $\mathrm{d}\Omega_{S^3} = \Omega^3\, \mathrm{d}^3p$, where the Jacobian $\Omega^3$ encodes the non-uniform compression of flat space onto the sphere.

\subsection{Inverse Projection}

The map is invertible away from the north pole:
\begin{equation}
p_i = \frac{p_0\, n_i}{1 - n_4}, \quad i = 1, 2, 3.
\label{eq:inverse}
\end{equation}
This confirms the projection is a diffeomorphism $S^3 \setminus \{N\} \to \mathbb{R}^3$, where $N = (0,0,0,1)$ is the north pole. The compactification of $\mathbb{R}^3$ to $S^3$ adjoins this single point at infinity, mirroring the Riemann sphere compactification $\mathbb{C} \cup \{\infty\} \cong S^2$ in one lower dimension.


%% ========================================================================
\section{Operator Equivalence and the Continuum Limit}
\label{sec:operator}
%% ========================================================================

\subsection{The Conformal Laplacian Identity}

Having established the geometric setting, we now demonstrate how the Laplace-Beltrami operator on $S^3$---the continuum limit of the graph Laplacian---relates to the flat-space Schr\"odinger operator.

For a conformal rescaling of the metric $g_{S^3} = \Omega^2 g_{\text{flat}}$ in $n = 3$ dimensions, the Laplace-Beltrami operators are related by
\begin{equation}
\Delta_{S^3} f = \Omega^{-2}\left[\Delta_{\text{flat}} f + \nabla(\ln\Omega) \cdot \nabla f\right],
\label{eq:conformal_laplacian}
\end{equation}
where $\Delta_{\text{flat}} = \partial^2/\partial p^2 + (2/p)\,\partial/\partial p$ is the radial Laplacian in three-dimensional flat space (for the $l = 0$ sector), and the connection correction $\nabla(\ln\Omega) \cdot \nabla f$ encodes the curvature of the stereographic coordinate system.

For the conformal factor $\Omega = 2p_0/(p^2 + p_0^2)$, the logarithmic gradient is
\begin{equation}
\frac{\mathrm{d} \ln\Omega}{\mathrm{d}p} = -\frac{2p}{p^2 + p_0^2}.
\label{eq:connection}
\end{equation}
This is the \textit{connection term} that distinguishes the curved Laplacian from the flat one. It is a rational function of $p$---no transcendental functions appear---which is essential for clean algebraic manipulation.

\subsection{Verification: Constants and Eigenfunctions}

As a consistency check, the Laplace-Beltrami operator annihilates constants:
\begin{equation}
\Delta_{S^3}(1) = 0,
\end{equation}
confirming the basic property that constants are harmonic on any Riemannian manifold.

For the first nontrivial eigenfunction, consider the hyperspherical angle $\chi$ defined by the stereographic projection $p = p_0 \tan(\chi/2)$. The function $\cos\chi$, which in stereographic coordinates takes the form
\begin{equation}
\cos\chi = \frac{p_0^2 - p^2}{p_0^2 + p^2},
\label{eq:cos_chi}
\end{equation}
is the $n = 2$, $l = 0$ eigenfunction of $\Delta_{S^3}$. Direct symbolic computation confirms
\begin{equation}
\Delta_{S^3}(\cos\chi) = -3\,\cos\chi.
\label{eq:eigenvalue_n2}
\end{equation}
More generally, the $l = 0$ eigenfunctions are Gegenbauer polynomials $C^1_{n-1}(\cos\chi)$:
\begin{align}
C^1_0 &= 1, \\
C^1_1 &= 2\cos\chi, \\
C^1_2 &= 4\cos^2\chi - 1,
\end{align}
with eigenvalues
\begin{equation}
\Delta_{S^3}\, C^1_{n-1}(\cos\chi) = -(n^2 - 1)\, C^1_{n-1}(\cos\chi).
\label{eq:eigenvalues}
\end{equation}
\textbf{[SYMBOLIC PROOF]} We have verified Eq.~\eqref{eq:eigenvalues} symbolically for $n = 1, 2, 3$ using the conformal Laplacian formula, confirming exact agreement.

\subsection{The Dimensionless Spectrum}

Equation~\eqref{eq:eigenvalues} is the central result of this section. The eigenvalues of the Laplace-Beltrami operator on the unit $S^3$ are
\begin{equation}
\lambda_n = -(n^2 - 1), \quad n = 1, 2, 3, \ldots
\label{eq:s3_spectrum}
\end{equation}
These are \textit{pure integers}. They carry no physical dimensions. They depend on no physical parameters---not on $p_0$, not on the electron mass, not on the fine structure constant. They are properties of the topology alone.

This spectrum arises from the same mechanism as the quantization of angular momentum on $S^2$: the compactness of the manifold forces the eigenfunctions to satisfy periodicity constraints that admit only discrete mode numbers. On $S^3$, these mode numbers correspond to the principal quantum number $n$ of hydrogen.

\subsection{The Energy Shell: Projecting to Physics}

How, then, do the dimensionful hydrogen energy levels emerge from dimensionless eigenvalues? The answer lies entirely in the energy-shell constraint Eq.~\eqref{eq:energy_shell}. Setting $p_0^2 = 1/n^2$ (hydrogen in atomic units), the eigenvalue equation on $S^3$ becomes an equation for the energy through the following algebraic chain:

\begin{enumerate}
\item The $S^3$ eigenvalue equation: $\Delta_{S^3}\,\Psi_n = -(n^2 - 1)\,\Psi_n$.
\item The energy-shell constraint: $p_0^2 = -2E = 1/n^2$.
\item Combine: the ratio $\lambda_n \cdot p_0^2 = -(n^2 - 1)/n^2 = -(1 - 1/n^2)$.
\item The physical energy: $E_n = -p_0^2/2 = -1/(2n^2)$.
\end{enumerate}

Step 4 is simply the definition of $p_0$ rearranged. The Rydberg formula is \textit{not} derived from the sphere's curvature. It is the energy-shell parameter, translated back into energy units. The integer eigenvalues $\lambda_n$ supply the quantum numbers; the energy shell supplies the units.

The complete transformation chain is summarized diagrammatically:
\begin{equation}
\boxed{
\begin{array}{c}
\text{Discrete Graph } \mathcal{G} \\[4pt]
\downarrow \;\text{\small continuum limit}\\[4pt]
\Delta_{S^3} \;\text{on unit } S^3 \\[4pt]
\lambda_n = -(n^2 - 1) \;\text{\small (dimensionless)}\\[4pt]
\downarrow \;\text{\small energy shell } p_0^2 = -2E\\[4pt]
E_n = -\frac{1}{2n^2} \;\text{\small (dimensionful)}
\end{array}
}
\label{eq:chain}
\end{equation}

\subsection{Conformal Decomposition of the Wavefunction}

To complete the operator equivalence, we demonstrate that the $S^3$ Laplacian acting on a conformally weighted wavefunction decomposes cleanly into the momentum-space Schr\"odinger structure.

Define the conformally weighted wavefunction
\begin{equation}
\Phi(p) = \Omega^2(p)\,\varphi(p),
\label{eq:conformal_wf}
\end{equation}
where $\varphi(p)$ is the physical momentum-space wavefunction. Applying $\Delta_{S^3}$ via the conformal identity Eq.~\eqref{eq:conformal_laplacian}, and expanding via the Leibniz rule, one obtains
\begin{equation}
\Delta_{S^3}[\Omega^2 \varphi] = \sum_{k=0}^{2} a_k(p)\, \frac{\mathrm{d}^k \varphi}{\mathrm{d}p^k},
\label{eq:decomposition}
\end{equation}
where the coefficients $a_k(p)$ are \textit{rational functions} of $p$ and $p_0$. No transcendental or unresolved terms appear. This decomposition has been verified symbolically: the coefficient of $\varphi''$ is nonzero (confirming a second-order kinetic operator exists), and all coefficients are manifestly rational, confirming that the $S^3$ Laplacian acting on conformally weighted wavefunctions produces exactly the structure of a second-order differential equation in the flat momentum coordinate.


%% ========================================================================
\section{Two-Electron Repulsion on $S^3$}
\label{sec:vee_s3}
%% ========================================================================

We now extend the $S^3$ framework to the two-electron repulsion
integral.

\subsection{Momentum-Space Representation}

The direct Coulomb integral (Slater $F^0$) between orbitals $a$ and $b$
in position space is
\begin{equation}
  F^0(a,b) = \int_0^\infty \!\!\int_0^\infty
  |R_a(r_1)|^2 \, |R_b(r_2)|^2 \, \frac{1}{r_>} \,
  r_1^2 \, r_2^2 \, dr_1 \, dr_2 \,,
  \label{eq:slater_f0}
\end{equation}
where $r_> = \max(r_1, r_2)$. For hydrogen-like orbitals with nuclear
charge $Z$, exact values are known:
\begin{align}
  F^0(1s,1s) &= \tfrac{5}{8}Z, \quad
  F^0(1s,2s) = \tfrac{17}{81}Z, \nonumber\\
  F^0(2s,2s) &= \tfrac{77}{512}Z.
  \label{eq:slater_exact}
\end{align}

An equivalent momentum-space representation exists for $s$-orbital pairs.
Defining the density Fourier transform
\begin{equation}
  \tilde{\rho}_a(q)
  = \frac{1}{q} \int_0^\infty |R_a(r)|^2 \, r \, \sin(qr) \, dr \,,
  \label{eq:density_ft}
\end{equation}
the Slater integral becomes
\begin{equation}
  F^0(a,b) = \frac{2}{\pi} \int_0^\infty
  \tilde{\rho}_a(q) \, \tilde{\rho}_b(q) \, dq \,.
  \label{eq:f0_momentum}
\end{equation}

For the hydrogen $1s$ and $2s$ orbitals:
\begin{align}
  \tilde{\rho}_{1s}(q) &= \frac{16Z^4}{(q^2 + 4Z^2)^2} \,,
  \label{eq:rho1s} \\[4pt]
  \tilde{\rho}_{2s}(q) &= \frac{Z^4(Z^4 - 3Z^2 q^2 + 2q^4)}{(Z^2 + q^2)^4} \,.
  \label{eq:rho2s}
\end{align}

\subsection{Fock Projection onto $S^3$}

We apply the Fock stereographic projection with parameter $p_0 = 2Z$,
under which $q = 2Z \tan u$ for $u \in [0, \pi/2)$ and
$dq = 2Z/\cos^2 u \, du$. Define the \emph{projected density}
\begin{equation}
  \Phi_a(t) \equiv \tilde{\rho}_a(2Zt) \,,
  \label{eq:projected_density}
\end{equation}
where $t = \tan u = q/(2Z)$ is the dimensionless momentum transfer.
Under this substitution, Eq.~\eqref{eq:f0_momentum} becomes
\begin{equation}
  \boxed{
    F^0(a,b) = \frac{4Z}{\pi} \int_0^\infty
    \Phi_a(t) \, \Phi_b(t) \, dt
  }
  \label{eq:f0_s3}
\end{equation}
\textbf{[SYMBOLIC PROOF]} This is the master formula: $F^0$ is a \emph{single} integral over one
$S^3$ coordinate, not a double integral over two points.

\subsubsection{$1s$--$1s$ verification}

Substituting $q = 2Zt$ into Eq.~\eqref{eq:rho1s}:
\begin{equation}
  \Phi_{1s}(t) = \frac{16Z^4}{(4Z^2 t^2 + 4Z^2)^2}
  = \frac{1}{(t^2 + 1)^2} \,.
  \label{eq:phi1s}
\end{equation}
Then
\begin{equation}
  F^0(1s,1s) = \frac{4Z}{\pi} \int_0^\infty
  \frac{dt}{(t^2+1)^4}
  = \frac{4Z}{\pi} \cdot \frac{5\pi}{32}
  = \frac{5Z}{8} \,. \quad\checkmark
\end{equation}

In angular coordinates ($t = \tan u$, $dt = du/\cos^2 u$),
the integrand becomes $\cos^6 u$, showing that the $1s$ density
is \emph{constant} on $S^3$ (up to the conformal Jacobian).

\subsubsection{$1s$--$2s$ verification (critical test)}

Substituting into Eq.~\eqref{eq:rho2s}:
\begin{equation}
  \Phi_{2s}(t) = \frac{1 - 12t^2 + 32t^4}{(4t^2 + 1)^4} \,.
  \label{eq:phi2s}
\end{equation}
This is \emph{not} constant on $S^3$---it has zeros at
$t^2 = (3 \pm 1)/16$, corresponding to nodes in the $2s$ radial
wavefunction. The integral evaluates to
\begin{align}
  F^0(1s,2s) &= \frac{4Z}{\pi} \int_0^\infty
  \frac{1 - 12t^2 + 32t^4}{(t^2+1)^2 (4t^2+1)^4} \, dt \nonumber\\
  &= \frac{4Z}{\pi} \cdot \frac{17\pi}{324}
  = \frac{17Z}{81} \,. \quad\checkmark
\end{align}

\subsubsection{$2s$--$2s$ verification}

\begin{align}
  F^0(2s,2s) &= \frac{4Z}{\pi} \int_0^\infty
  \frac{(1 - 12t^2 + 32t^4)^2}{(4t^2+1)^8} \, dt \nonumber\\
  &= \frac{4Z}{\pi} \cdot \frac{77\pi}{2048}
  = \frac{77Z}{512} \,. \quad\checkmark
\end{align}

\textbf{[SYMBOLIC PROOF]} All three Slater integrals are reproduced exactly by the $S^3$
density-overlap formula~\eqref{eq:f0_s3} (verified symbolically in
\texttt{tests/test\_paper7\_vee\_s3.py}: the three integrals evaluate to
$5\pi/32$, $17\pi/324$, $77\pi/2048$ in exact arithmetic, reproducing
\eqref{eq:slater_exact}).

\subsection{Structural Result: Node Property vs.\ Edge Property}
\label{sec:node_vs_edge}

The key structural finding is that
Eq.~\eqref{eq:f0_s3} is a \emph{single} integral over one $S^3$
coordinate (the momentum transfer $t$), not a double integral over
two independent $S^3$ points. Each orbital $a$ contributes a density
\emph{function} $\Phi_a(t)$ on $S^3$, and $V_{ee}$ is the weighted
overlap:
\begin{equation}
  V_{ee}(a,b) = \langle \Phi_a | \Phi_b \rangle_{S^3}
  \equiv \frac{4Z}{\pi} \int_0^\infty \Phi_a(t) \, \Phi_b(t) \, dt \,.
\end{equation}

In graph-theoretic language, $V_{ee}$ is a \textbf{node property}
(each node carries a density profile) rather than an \textbf{edge
property} (a pairwise distance between node coordinates).

\paragraph{Failure of the chordal ansatz.}
A naive chordal-distance ansatz $V_{ee} = \kappa_{ee}/d^2_{\text{chord}}$,
where $d^2 = 2(1 - \boldsymbol{\xi}_a \cdot \boldsymbol{\xi}_b)$
is the chordal distance between $S^3$ points assigned to each orbital,
gives (with $\kappa_{ee} = 5Z/2$):
\begin{center}
\begin{tabular}{lccc}
\hline
Pair & Slater $F^0$ & $\kappa/d^2$ & Ratio \\
\hline
$1s$--$1s$ & $5Z/8$ & $5Z/8$ & 1.00 \\
$1s$--$2s$ & $17Z/81$ & $25Z/4$ & 29.8 \\
$2s$--$2s$ & $77Z/512$ & $5Z/8$ & 4.2 \\
\hline
\end{tabular}
\end{center}
The chordal ansatz is calibrated to reproduce $F^0(1s,1s)$ exactly
(because the $1s$ density is constant on $S^3$, making a single point
adequate), but overestimates $F^0(1s,2s)$ by a factor of~30.
The fundamental reason is that the $2s$ density function has nodes
and oscillations on $S^3$ that cannot be captured by a single coordinate
point.

\subsection{Natural Projection Scales}

Each shell has a natural Fock projection scale:
\begin{equation}
  n = 1: \quad p_0 = 2Z \qquad\text{(denominator } q^2 + 4Z^2\text{)} \,,
\end{equation}
\begin{equation}
  n = 2: \quad p_0 = Z \qquad\text{(denominator } q^2 + Z^2\text{)} \,.
\end{equation}
When all orbitals are projected using a \emph{single} $p_0 = 2Z$, the
$n=1$ density becomes the constant function $\Phi_{1s} = 1/(1+t^2)^2$,
while higher-shell densities acquire polynomial modulations. These
modulations encode the radial node structure and are responsible for
the reduced overlap integrals ($17/81 < 5/8$).

\subsection{Implications}
\label{sec:vee_implications}

The density-overlap formula~\eqref{eq:f0_s3} has three implications
for the computational framework:

\begin{enumerate}
  \item \textbf{Exact Slater integrals on $S^3$.}
    The Fock projection converts position-space double radial integrals
    into single rational-function integrals on $S^3$, which can be
    evaluated in closed form via partial fractions.

  \item \textbf{Graph implementation.}
    In the graph Laplacian framework, each node (orbital state) should
    carry a density profile $\Phi_a(t)$, and $V_{ee}$ matrix elements
    should be computed as overlap integrals rather than pairwise edge
    weights.

  \item \textbf{Angular momentum extension.}
    The $1D$ monopole formula~\eqref{eq:f0_momentum} applies to
    $s$-orbital pairs. For $l > 0$, the full $3D$ angular convolution
    is required, involving the multipole expansion of the orbital
    density in spherical harmonics.
\end{enumerate}


%% ========================================================================
\section{Generalization to $N$ Electrons}
\label{sec:N_electron}
%% ========================================================================

The preceding sections established the $S^3$ structure for a single
electron.  The framework generalizes naturally to $N$-electron atoms:
the $N$-body configuration space $\mathbb{R}^{3N}$ maps, via the
hyperspherical Fock projection, onto the unit sphere
$S^{3N-1}$~\cite{loutey_paper13}.

\subsection{The Angular Manifold and Its Isometry Group}

For $N$ electrons at positions $\bm{r}_1, \ldots, \bm{r}_N \in
\mathbb{R}^3$, the hyperradius $R = \sqrt{\sum_i r_i^2}$ separates
a radial coordinate from a $(3N-1)$-dimensional angular manifold.
In the limit $R \to 0$ (all electrons at the nucleus), the angular
dynamics reduces to the Laplace--Beltrami operator on the unit
$S^{3N-1}$, whose isometry group is $\mathrm{SO}(3N)$.

The eigenvalues of the grand angular momentum operator $\Lambda^2$
on $S^{3N-1}$ are the $\mathrm{SO}(3N)$ Casimir eigenvalues:
\begin{equation}
  \mu_{\mathrm{free}}(\nu) = \frac{\nu(\nu + 3N - 2)}{2},
  \quad \nu = 0, 1, 2, \ldots
  \label{eq:casimir_general}
\end{equation}
These are the ``free'' eigenvalues---the spectrum before any
electron-electron or electron-nuclear interactions are included
in the charge function $C/R$.

\subsection{The Hierarchy of Spheres}

Table~\ref{tab:N_electron} summarizes the angular manifold, symmetry
group, and Casimir structure for the first few electron numbers.
Each row represents a qualitatively different angular problem,
but all share the same structural origin: the Laplace--Beltrami
operator on a unit sphere of appropriate dimension.

\begin{table}[h]
\caption{Angular manifold structure for $N$-electron atoms.
The centrifugal term $l_{\mathrm{eff}}(l_{\mathrm{eff}}+1)/(2R^2)$
arises from extracting the Jacobian factor $R^{-(3N-1)/2}$.
The effective quantum number $n_{\mathrm{eff}} = l_{\mathrm{eff}} + 1$
determines the quasi-Coulomb energy.}
\label{tab:N_electron}
\begin{ruledtabular}
\begin{tabular}{lccccc}
  $N$ & $S^{d-1}$ & $\mathrm{SO}(d)$ & $\mu_{\mathrm{free}}$
  & $l_{\mathrm{eff}}$ & $n_{\mathrm{eff}}$ \\
  \hline
  1 & $S^2$ & SO(3) & $\nu(\nu+1)/2$   & 0   & 1 \\
  2 & $S^5$ & SO(6) & $\nu(\nu+4)/2$   & 3/2 & 5/2 \\
  3 & $S^8$ & SO(9) & $\nu(\nu+7)/2$   & 3   & 4 \\
  4 & $S^{11}$ & SO(12) & $\nu(\nu+10)/2$ & 9/2 & 11/2 \\
\end{tabular}
\end{ruledtabular}
\end{table}

For hydrogen ($N = 1$), $S^2$ recovers the familiar angular momentum
eigenvalues $l(l+1)/2$ and the $\mathrm{SO}(3)$ rotation group, while
$S^3$ (from the full 3D problem including the radial Fock projection)
gives the $\mathrm{SO}(4)$ hydrogen symmetry of Section~\ref{sec:conformal}.
For helium ($N = 2$), the angular manifold is $S^5$ with
$\mathrm{SO}(6)$ symmetry, and the free eigenvalues
$\nu(\nu+4)/2 = 0, 6, 16, 30, \ldots$ (even-$\nu$ singlet sector)
have been verified numerically
to $< 0.001\%$ against the hyperspherical angular
solver~\cite{loutey_paper13}.

\subsection{Fermionic Statistics and Ground-State Selection}

The permutation group $S_N$ acts on the angular manifold by
exchanging electron coordinates.  Fermionic antisymmetry requires
the spatial wavefunction to transform under a specific irreducible
representation (irrep) of $S_N$, determined by the spin state.

For helium ($N = 2$, singlet ${}^1S$), the spatial part is
symmetric under $S_2$ exchange $\alpha \to \pi/2 - \alpha$.
This selects even grand angular momentum $\nu = 0, 2, 4, \ldots$,
and the ground state sits at $\nu = 0$ with $\mu_{\mathrm{free}} = 0$.

For lithium ($N = 3$, doublet ${}^2S$), the spatial part transforms
as the standard (mixed-symmetry) irrep $[2,1]$ of $S_3$.  This
irrep first appears at $\nu = 1$, so the lowest allowed free
eigenvalue is $\mu_{\mathrm{free}} = 1 \cdot 8/2 = 4$---not zero.
The higher baseline reflects the Pauli exclusion principle encoded
group-theoretically: more electrons require higher angular momentum
to satisfy antisymmetry.

\subsection{Implications for the Dimensionless Vacuum}

The $N$-electron generalization preserves the core principle: at
$R = 0$, the angular dynamics is a free particle on the unit
$S^{3N-1}$, with pure integer-ratio eigenvalues determined by
the $\mathrm{SO}(3N)$ Casimir.  No physical scale enters.  The
electron-nuclear and electron-electron interactions enter through
the charge function $C(\hat{\Omega})/R$, which acts as a
perturbation that lifts the Casimir degeneracy---analogous to how
the energy-shell constraint $p_0^2 = -2E$ projects the dimensionless
$S^3$ eigenvalues into dimensionful energies for the single-electron
case.

The 18 symbolic proofs of Appendix~\ref{app:tests} verify the
$N = 1$ case.  \textbf{[CONJECTURE]} The generalization to $N > 1$ is expected by the same
algebraic structure: the Laplace--Beltrami operator on $S^{3N-1}$
is the continuum limit of a discrete graph Laplacian on the
$N$-electron angular lattice, and its spectrum is determined by
the topology of the sphere alone.


%% ========================================================================
\section{Discussion and Conclusions}
\label{sec:discussion}
%% ========================================================================

\subsection{Summary of Results}

We have established a mathematical chain connecting the discrete GeoVac lattice to the continuous Schr\"odinger equation through three stages---the first empirical, the second and third exact and algebraic:

\textit{Stage 1: Discrete $\to$ Continuous.} The paraboloid lattice, with $n^2$-dimensional shells connected by $SU(2) \otimes SU(1,1)$ transition operators, is read as converging in the continuum limit to the Laplace-Beltrami operator on a compact manifold with $SO(4)$ isometry (measured: the saturation $\lambda_{\max} \to 2 d_{\max}$ and the decay of the $s/p$ lift; the operator convergence itself is not established, New Contributions item 1).

\textit{Stage 2: Curved $\to$ Flat.} Fock's stereographic projection maps the unit $S^3$ to flat $\mathbb{R}^3$ momentum space, with conformal factor $\Omega = 2p_0/(p^2 + p_0^2)$. The chordal distance identity (Eq.~\ref{eq:chordal}) converts the free-particle Green's function on $S^3$ into the Coulomb kernel in flat space.

\textit{Stage 3: Dimensionless $\to$ Dimensionful.} The energy-shell constraint $p_0^2 = -2E$ converts pure integer eigenvalues $\lambda_n = -(n^2 - 1)$ into the Rydberg energy levels $E_n = -1/(2n^2)$.

Stages 2 and 3 are verified by independent symbolic computation: the 18 passing tests constitute an algebraic audit of the conformal transformation chain (ten of the eighteen are definitional, limiting-case or non-discriminating checks; the load-bearing identities are items 1, 6, 7, 8, 10, 11, 16 and 17), covering the unit sphere constraint, the conformal factor identities, the projection limits, the chordal distance identity, the dot-product form, the volume Jacobian, the inverse projection, the connection term, the conformal Laplacian of $\Omega^2$, the zero-mode annihilation, the eigenvalue equation for $n = 2$ and $n = 3$, the conformal decomposition, and the eigenvalue-to-energy mapping.  Stage 1---the graph-to-continuum convergence itself---is only indicated numerically through $n_{\max} = 30$ (spectral-bound saturation and $s/p$-lift decay; the operator convergence itself is not established) with no rate proven here (see Limitations; the operator-system rate results of Paper~38~\cite{loutey_paper38} concern the truncated spectral triple, whose identification with this graph Laplacian is not established in this paper).

\subsection{The Scale-Invariance Property}

As Fock showed, the Schr\"odinger equation for hydrogen is mathematically equivalent, under conformal projection, to the Laplace-Beltrami operator on the unit $S^3$---a dimensionless, scale-invariant topology. What we add is the observation that the GeoVac discrete graph converges to this same structure, so the continuous equation and its dimensionful energy levels can be viewed as emergent properties of projecting a discrete topology into observational coordinates. This is a mathematical equivalence, not a proof of physical priority---whether the graph or the differential equation is ``more fundamental'' is a question of physical interpretation, not of theorem.

Three features of the projection merit emphasis:

\textit{(i) Scale invariance.} The conformal factor $\Omega$ always produces the unit $S^3$, regardless of $p_0$. Substituting $u = p/p_0$ yields the standard metric $4/(1+u^2)^2\,\mathrm{d}u^2$ with all $p_0$-dependence cancelled. The intrinsic geometry of the vacuum has no scale.

\textit{(ii) Integer spectrum.} The eigenvalues $\lambda_n = -(n^2-1)$ are pure integers determined by the topology of $S^3$. They count the number of independent oscillation modes that fit on the compact manifold, in exact analogy with standing waves on a finite string.

\textit{(iii) Energy as projection cost.} The dimensionful energy $E_n = -1/(2n^2)$ enters exclusively through the energy-shell constraint---the requirement that the stereographic focal length be consistent with the bound-state energy. Energy is the cost of representing curved topological information in flat coordinates.

\subsection{Relation to the Coulomb Potential}

A frequently asked question is: where does the $1/r$ Coulomb potential come from? In the standard formulation, $V(r) = -1/r$ is an input to the Schr\"odinger equation. As Fock showed~\cite{fock1935}, in the $S^3$ formulation it is an \textit{output} of the conformal projection.

The mechanism is the chordal distance identity, Eq.~\eqref{eq:chordal}. On $S^3$, the Green's function for the Laplace-Beltrami operator is $G(\mathbf{n}, \mathbf{n}') \propto 1/|\mathbf{n} - \mathbf{n}'|^2$, which is the natural interaction kernel on a sphere. Under stereographic projection, this becomes
\begin{equation}
\frac{1}{|\mathbf{n} - \mathbf{n}'|^2} = \frac{1}{\Omega(\mathbf{p})\,\Omega(\mathbf{p}')\,|\mathbf{p} - \mathbf{p}'|^2},
\end{equation}
which, when Fourier transformed to position space, yields the $1/r$ Coulomb potential (with conformal weight corrections that are absorbed into the wavefunction normalization). The Coulomb singularity at $r = 0$ corresponds to the divergence of the stereographic map at the south pole---it is a coordinate singularity of the projection, not a physical singularity of the underlying topology.

\subsection{Implications for the Discrete Model}

The results of this paper provide a formal justification for the GeoVac framework's computational approach. Reading the continuum limit of the graph Laplacian as conformally equivalent to the unit $S^3$ (an Observation; New Contributions, item 1) accounts for several empirical observations:

\textbf{[PANEL-VERIFIED]} \textit{Spectral exactness:} The algebraic operators $T_\pm, L_\pm$ reproduce the exact Rydberg \emph{quantum numbers} $n$ and $\ell(\ell+1)$ because they encode the $SO(4)$ Casimir structure of $S^3$; the energies follow only after the matched scale $\kappa$ is applied, which is an observation and not a derivation (Section~\ref{sec:discussion}).

\textbf{[INTERNAL THEOREM]} \textit{Discretization artifacts:} The $s/p$ node-amplitude proxy at finite $n_{\max}$ has a closed-form value on $n_{\max} \equiv 0 \ (\mathrm{mod}\ 3)$ and decays as $n_{\max} \to \infty$, but its size is set by that residue rather than by boundary effects;\ the ``spectral aliasing'' account is withdrawn (2026-09-03), the $\ell$-blocks being disconnected.

\textit{Kinetic scale factor:} The constant $\kappa = -1/16$ is the
kinetic scale factor that bridges the dimensionless graph and
dimensionful physics.  Its value $1/16$ coincides with a
conformal-geometry quantity, as follows.  The $l = 0$ eigenfunctions $C^1_{n-1}(\cos\chi)$ of
Eq.~\eqref{eq:eigenvalues} are the Chebyshev polynomials of the second
kind $U_{n-1}(\cos\chi)$, whose three-term recurrence $\cos\chi\,U_{n-1}
= \tfrac12 U_n + \tfrac12 U_{n-2}$ gives a normalized transition
amplitude of $1/2$ between adjacent shells.  The squared inter-shell
coupling of $\cos\chi$ in the Gegenbauer eigenbasis is
\textbf{[SYMBOLIC + MEASURED]}
\begin{equation}
  c^2(n,l) \;:=\; \bigl|\langle n{+}1, l\,|\cos\chi|\,n, l\rangle\bigr|^2
  \;=\; \frac{1}{4}\!\left[ 1 - \frac{l(l+1)}{n(n+1)} \right],
  \label{eq:fock_coupling}
\end{equation}
(closed form from the Gegenbauer recurrence and norm ratio; verified by
quadrature at seven $(n,l)$ pairs,
\texttt{tests/test\_paper7\_gegenbauer\_coupling.py}).  For $l = 0$,
$c^2(n,0) = 1/4$ for all~$n$.  \emph{Correction (2026-09-03):} an
earlier version printed this formula with prefactor $1/16$ and read it
as ``$(1/4)^2$''; the amplitude is $1/2$, so the coupling is $1/4$, and
the $1/16$ is a \emph{different} quantity.  The geometric quantity that
the matching $\kappa$ coincides with is
\begin{enumerate}
  \item the inverse Fock Jacobian $1/\Omega^4(0) = 1/16$, where $\Omega(0) = 2$
        is the stereographic conformal factor at the projection center
        ($p = 0$) in the dimensionless $p_0 = 1$ convention
        (dimensionfully $\Omega(0) = 2/p_0$; the $\kappa$
        observation is stated in the dimensionless convention);
  \item equivalently $c^2(n,0)/4$, one quarter of the $l = 0$ base rate of
        the Casimir decomposition in Eq.~\eqref{eq:fock_coupling}---an
        arithmetic restatement of item 1, not an independent reading.
\end{enumerate}
\textbf{[OBSERVATION]} We record this as an \emph{observation}, not a derivation:\ the
energy-matching $\kappa = -1/16$ (the scale mapping the graph's
$\lambda_{\max} \to 8$ onto $-\tfrac12$) coincides numerically with the
geometric $1/16$ above, but the two are computed from disjoint
inputs---the energy target and graph coordination on the matching
side, the conformal factor on the geometric side---and no mechanism is
known that carries one onto the other (a counterfactual energy target
moves the matching $\kappa$ while the geometric $1/16$ is rigid;
falsifier \texttt{tests/test\_trunk\_qa\_kappa.py}).  The universality
of $\kappa$ across nuclear charges $Z$ (via $Z^2$ scaling) confirms
that the graph topology correctly encodes the \textit{relative}
spectral structure.

A notable special value: at $(n,l) = (4,3)$ the Casimir factor is
$1 - 12/20 = 2/5$, and $(2/5) \times 1/\Omega^4(0) = 1/40 = \Delta$, the
boundary term in Paper~2's Observation-tier combination rule
$K = \pi(B + F - \Delta)$~\cite{loutey_paper2}.  The value $\Delta =
1/40$ is independently derived from the Dirac degeneracy
$g_3^{\mathrm{Dirac}} = 40$ identified in Phase~4H; the composite
$[1 - l(l+1)/(n(n+1))]/\Omega^4(0)$ is not a matrix element of
anything (the actual coupling is $c^2(4,3) = 1/10$),
\textbf{[OBSERVATION]} so the coincidence $\Delta = (2/5)/\Omega^4(0)$
is a numerical observation about a composite of two geometric
quantities, not a two-route confirmation
(\texttt{tests/test\_trunk\_qa\_c2\_delta.py} records exactly this
status; before 2026-09-03 it was printed as $c^2(4,3) = 1/40$ under the
$1/16$-prefactor formula).  The coupling itself is now derived
(\texttt{tests/test\_paper7\_gegenbauer\_coupling.py}); the former
``named upgrade path'' is closed.

\subsection{Limitations}

Several caveats apply. First, the single-electron framework of Sections~\ref{sec:graph}--\ref{sec:operator} addresses hydrogen only. Section~\ref{sec:vee_s3} extends the framework to two-electron repulsion for $s$-orbital pairs; the general case with angular momentum requires the full three-dimensional Fourier convolution. Second, the conformal projection presented here recovers only the non-relativistic Schr\"odinger equation. Relativistic corrections (fine structure, spin-orbit coupling) presumably require additional geometric structure, such as the fiber bundle construction explored in companion papers \cite{loutey_paper0,loutey_paper1}. \textbf{[MEASURED]} Third, while the discretization artifacts of the graph Laplacian decay systematically, we have not proven a rate of convergence here---only demonstrated it empirically through $n_{\max} = 30$.  The rigorous state-space Gromov--Hausdorff convergence of the scalar truncations for compact metric groups (including $\mathrm{SU}(2) = S^3$) is established by Gaudillot-Estrada and van Suijlekom~\cite{gaudillot_vs2025}; the Dirac/spinor analog, with explicit rate $(4/\pi)\log n/n$, is proved in Paper~38 of this series.

\subsection{Concluding Remarks}

Vladimir Fock discovered in 1935 that the hydrogen atom, when viewed in momentum space, lives naturally on a three-sphere. The scale invariance of this sphere---the fact that it is intrinsically the \textit{unit} $S^3$ regardless of the energy-shell parameter---was implicit in Fock's treatment but has not, to our knowledge, been stated as an explicit organizing principle. We have articulated this scale-invariance property and, more importantly, \textbf{[OBSERVATION]} recorded the numerical indications, through $n_{\max} = 30$, that the GeoVac discrete graph Laplacian converges to the same $S^3$ structure in the continuum limit (the spectral-bound saturation $\lambda_{\max} \to 2 d_{\max}$ and the $s/p$-lift decay; the operator convergence itself is not established, 2026-09-03).  That reading is the empirical foundation of the sparse graph computational approach; the conformal chain from the limit onward is what is proven symbolically.

Within this mathematical equivalence, the wave equation can be viewed as a continuous, dimensionful projection of a discrete, dimensionless topology. The Coulomb potential emerges as the coordinate distortion created by the stereographic flattening, and the energy levels emerge from the energy-shell constraint that parameterizes the projection.

We emphasize that these are statements of mathematical equivalence, not claims about which formulation is physically prior. Whether the topological perspective points toward a deeper theory of quantum mechanics, or provides a computationally useful reformulation, is a question we leave to future investigation. \textbf{[SYMBOLIC PROOF]} What is established here is the precise mathematical chain connecting the graph's \emph{continuum limit} to the continuous equation, verified step by step through symbolic computation; the convergence of the graph to that limit is neither proven nor established numerically---it is indicated by the spectral-bound saturation alone --- the $s/p$ leg was withdrawn as evidence on 2026-09-03, its size being set by the residue of $n_{\max}$ modulo $3$ --- and recorded as an Observation.

The $S^3$ conformal equivalence established here is one instance of a broader pattern: the discrete graph topology defined by the quantum number node structure $\{(n, l, m)\}$ admits multiple conformally distinct continuum embeddings, of which $S^3$ is the natural one for the single-center Coulomb problem.  The same combinatorial skeleton reappears in prolate spheroidal coordinates for one-electron diatomics~\cite{loutey_paper11}, hyperspherical coordinates for two-electron atoms~\cite{loutey_paper13}, molecule-frame hyperspherical coordinates for two-electron molecules, and composed fiber bundles for core--valence systems.  The full hierarchy of these embeddings, and the sense in which the discrete node structure serves as a combinatorial invariant across them, is catalogued in Paper~0~\cite{loutey_paper0}.

\subsection{Extension to Molecules}

The principle that the graph Laplacian should discretize the Laplace-Beltrami operator on the natural conformal manifold extends beyond single-center systems.  For diatomic molecules, the two-center Coulomb problem separates in prolate spheroidal coordinates $(\xi, \eta, \varphi)$, just as the single-center problem separates in spherical coordinates (leading to $S^3$ via Fock's projection).  Discretizing the Laplace-Beltrami operator on the prolate spheroid yields a molecular lattice with $R$-dependence built into the metric---no fitted parameters are required.  \textbf{[MEASURED]} For H$_2^+$, this approach achieves $R_{\mathrm{eq}} = 2.001$~bohr (0.21\% error) and $E_{\min} = -0.598$~Ha (0.70\% error) with zero free parameters~\cite{loutey_paper11}.  \textbf{[CONJECTURE]} This suggests a general principle: every separable quantum system has a natural lattice defined by the coordinate system in which separation occurs.


%% ========================================================================
% References
%% ========================================================================

\begin{thebibliography}{99}

\bibitem{fock1935}
V. Fock,
``Zur Theorie des Wasserstoffatoms,''
Z. Phys. \textbf{98}, 145--154 (1935).

\bibitem{bargmann1936}
V. Bargmann,
``Zur Theorie des Wasserstoffatoms: Bemerkungen zur gleichnamigen Arbeit von V. Fock,''
Z. Phys. \textbf{99}, 576--582 (1936).

\bibitem{barut1967}
A.~O. Barut and H. Kleinert,
``Transition probabilities of the hydrogen atom from noncompact dynamical groups,''
Phys. Rev. \textbf{156}, 1541--1545 (1967).

\bibitem{barut1967b}
A.~O. Barut and H. Kleinert,
``Current operators and Majorana equation for the hydrogen atom from dynamical groups,''
Phys. Rev. \textbf{157}, 1180 (1967).

\bibitem{bander1966}
M. Bander and C. Itzykson,
``Group theory and the hydrogen atom (I),''
Rev. Mod. Phys. \textbf{38}, 330--345 (1966).

\bibitem{biedenharn1981}
L.~C. Biedenharn and J.~D. Louck,
\textit{Angular Momentum in Quantum Physics},
Encyclopedia of Mathematics and its Applications, Vol. 8 (Addison-Wesley, Reading, MA, 1981).

\bibitem{loutey_paper0}
J. Loutey,
``Angular Momentum as Information Geometry: Isotropic Packing and the
Universal Angular Cross-Section,''
GeoVac Paper 0 (2026).

\bibitem{loutey_paper1}
J. Loutey,
``The Geometric Atom: Atomic Structure as a Packing Problem,''
GeoVac Paper 1 (2026).

\bibitem{loutey_paper2}
J. Loutey,
``The Fine Structure Constant as a Spectral Coincidence on the Hopf
Fibration,''
GeoVac Paper 2 (2026).

\bibitem{loutey_paper11}
J. Loutey,
``The Molecular Fock Projection: Diatomic Molecules as Prolate
Spheroidal Lattices,''
GeoVac Paper 11 (2026).

\bibitem{loutey_paper13}
J. Loutey,
``The Hyperspherical Lattice: Two-Electron Atoms as Coupled Channel
Graphs,''
GeoVac Paper 13 (2026).

\bibitem{loutey_paper18}
J.~Loutey,
``Spectral-Geometric Exchange Constants: Projecting Discrete Spectra onto Continuous Manifolds,''
GeoVac Paper~18 (2026).
\bibitem{loutey_paper38}
J.~Loutey,
``State-space Gromov--Hausdorff convergence of truncated Camporesi--Higuchi spectral triples on $S^3$,''
GeoVac Paper~38 (2026).
\bibitem{gaudillot_vs2025}
Y.~Gaudillot-Estrada and W.~D.~van~Suijlekom,
``Convergence of spectral truncations for compact metric groups,''
\textit{Int.~Math.~Res.~Not.\ IMRN} \textbf{2025}, no.~13, rnaf197 (2025); arXiv:2310.14733.

\end{thebibliography}


%% ========================================================================
\appendix
%% ========================================================================

\section{Symbolic Test Suite Summary}
\label{app:tests}

The following 18 tests were implemented in \texttt{sympy} (Python symbolic algebra) and executed with all assertions passing. Each test verifies a specific algebraic identity in the conformal transformation chain. The tests are organized into two modules.

\subsection{Module I: Stereographic Projection Geometry (\texttt{tests/test\_fock\_projection.py})}

\begin{enumerate}
\item \textbf{Unit sphere constraint.} $\sum_i n_i^2 = 1$ for all $\mathbf{p} \in \mathbb{R}^3$.
\item \textbf{Conformal factor identity.} $n_1^2 + n_2^2 + n_3^2 = \Omega^2 |\mathbf{p}|^2$.
\item \textbf{South pole limit.} $\mathbf{p} = 0 \Rightarrow \mathbf{n} = (0,0,0,-1)$.
\item \textbf{North pole limit.} $|\mathbf{p}| \to \infty \Rightarrow \mathbf{n} \to (0,0,0,+1)$.
\item \textbf{Conformal factor limits.} $\Omega(0) = 2/p_0$; $\Omega(\infty) = 0$.
\item \textbf{Chordal distance identity.} $|\mathbf{n} - \mathbf{n}'|^2 = \Omega\,\Omega'\,|\mathbf{p} - \mathbf{p}'|^2$.
\item \textbf{Dot product form.} $\mathbf{n} \cdot \mathbf{n}' = 1 - \tfrac{1}{2}\Omega\,\Omega'\,|\mathbf{p} - \mathbf{p}'|^2$.
\item \textbf{Volume Jacobian.} The induced metric of the embedding is $g = J^{\top}J = \Omega^2\, I_3$ (conformality, all nine entries), hence $\sqrt{\det g} = \Omega^3$; computed from the $4 \times 3$ Jacobian of the embedding, not restated.
\item \textbf{Numerical spot-check.} All identities verified at concrete floating-point values.
\item \textbf{Inverse projection.} $p_i = p_0\, n_i/(1 - n_4)$ round-trips to the identity.
\end{enumerate}

\subsection{Module II: Conformal Laplacian and Eigenvalues (\texttt{tests/test\_fock\_laplacian.py})}

\begin{enumerate}\setcounter{enumi}{10}
\item \textbf{Connection term.} $\mathrm{d}(\ln\Omega)/\mathrm{d}p = -2p/(p^2 + p_0^2)$.
\item \textbf{Laplacian of $\Omega^2$.} $\Delta_{\text{flat}}(\Omega^2)/\Omega^2$ is a rational function.
\item \textbf{Zero-mode annihilation.} $\Delta_{S^3}(1) = 0$.
\item \textbf{$n=2$ eigenvalue.} $\Delta_{S^3}(\cos\chi) = -3\cos\chi$.
\item \textbf{Conformal decomposition.} $\Delta_{S^3}[\Omega^2 \varphi]$ has rational coefficients in $\{\varphi, \varphi', \varphi''\}$.
\item \textbf{$n=2$ eigenfunction.} Full eigenvalue equation verified symbolically.
\item \textbf{$n=3$ eigenfunction.} $\Delta_{S^3}(4\cos^2\chi - 1) = -8(4\cos^2\chi - 1)$.
\item \textbf{Eigenvalue-to-energy mapping.} $\lambda_n = -(n^2-1)$ with $p_0^2 = 1/n^2$ gives $E_n = -1/(2n^2)$ (definitional: the energy-shell relation rearranged, not an independent identity).
\end{enumerate}


\section{Conformal Laplacian Derivation}
\label{app:derivation}

For completeness, we derive Eq.~\eqref{eq:conformal_laplacian}. Let $\tilde{g}_{ij} = \Omega^2 g_{ij}$ be a conformally rescaled metric in $n$ dimensions, with $\Omega = e^\omega$. The Christoffel symbols transform as
\begin{equation}
\tilde{\Gamma}^k_{ij} = \Gamma^k_{ij} + \delta^k_i\, \partial_j \omega + \delta^k_j\, \partial_i \omega - g_{ij}\, g^{kl}\, \partial_l \omega.
\end{equation}
The Laplace-Beltrami operator $\Delta_{\tilde{g}} f = \tilde{g}^{ij}\left(\partial_i \partial_j f - \tilde{\Gamma}^k_{ij}\, \partial_k f\right)$ yields, after substitution:
\begin{equation}
\Delta_{\tilde{g}} f = \Omega^{-2}\left[\Delta_g f + (n-2)\, g^{ij}\, (\partial_i \omega)(\partial_j f)\right].
\end{equation}
Setting $n = 3$ and $g = g_{\text{flat}}$ gives Eq.~\eqref{eq:conformal_laplacian}, with the connection correction proportional to $(n - 2) = 1$.

\paragraph{Note added v2.6.0.}
The $S^3$ density overlap formula (Sec.~V) provides exact rational
Slater integrals for all hydrogenic orbital pairs, enabling a fully
graph-native FCI calculation for helium.  Using a \emph{hybrid} one-body operator --- the exact hydrogenic
eigenvalues $-Z^2/2n^2$ on the diagonal and the graph Laplacian's
off-diagonal entries scaled by $\kappa$, the degree term dropped
(\texttt{geovac/casimir\_ci.py}) --- together with rational Slater integrals as
the two-body operator, the He ground state is computed with \emph{no fitted
parameters}: no orbital exponent, no variational optimization, no numerical
quadrature.  \emph{Corrected 2026-09-03:} this passage read ``using the graph
Laplacian $\kappa\mathcal{L}$ as the one-body operator \ldots\ with zero free
parameters beyond the universal matched scale $\kappa$'', which overstates what
is computed.  The exact one-electron spectrum is an \emph{input} here rather
than a graph output --- $\kappa\mathcal{L}$ alone has $\lambda_{\max} = 0$ at
$n_{\max} = 1$ and would give $E = +5Z/8$, not the value below --- so this
calculation is not a demonstration that the graph supplies the one-body
physics.  The FCI matrix at each
$n_{\max}$ is a rational matrix whose ground-state eigenvalue is an
algebraic number.

At $n_{\max} = 7$ (1,218 singlet configurations), the graph-native
\textbf{[MEASURED]} energy is $-2.8983$~Ha (0.19\% error vs.\ exact $-2.903724$~Ha),
\textbf{[OBSERVATION]} with a convergence floor estimated at $\sim$0.14\%.
The graph's off-diagonal one-body structure (inter-shell $T_\pm$
couplings) is the dominant ingredient: replacing the graph Laplacian
\textbf{[MEASURED]} with the diagonal hydrogenic operator degrades the result to 1.41\%
(variational exponent) or 1.85\% (fixed $k = Z$).

FCI basis invariance was verified computationally (Sprint~3D):
transforming $V_{ee}$ to the graph eigenbasis yields energies
identical to the graph-native result (to machine precision,
$< 3 \times 10^{-15}$~Ha), \textbf{[OBSERVATION]} consistent with the convergence floor
arising from the electron-electron cusp rather than from basis mismatch
(embedding exchange constant, Paper~18~\cite{loutey_paper18}).

At $n_{\max} = 1$, the graph-native energy reduces to
\textbf{[SYMBOLIC PROOF]} $E = -Z^2 + 5Z/8 = -11/4$~Ha for $Z = 2$, an exact rational
(two one-body terms $-Z^2/2$ plus the repulsion $F^0(1s,1s) = 5Z/8$;
an earlier printing halved the one-body sum and gave $-11/8$).
\textbf{[SYMBOLIC PROOF]} The variationally optimized single-exponent result $k^* = 27/16$,
$E^* = -729/256$~Ha is also an exact rational.
A 2D variational solver treating hyperradius and hyperangle
simultaneously (Paper~13, Note added v2.6.0) achieves 0.004\% with
cusp correction; both methods sit above the exact value, so they
do not bracket it.

\end{document}
